Los códigos QR se volvieron parte de la rutina diaria mucho más rápido de lo que
avanzaron las medidas de seguridad que deberían acompañarlos. Cuando alguien escanea
un código, no ve a dónde lo lleva el enlace antes de que su teléfono lo abra, y la
mayoría de sistemas operativos móviles hacen esa redirección sin revisar nada. Ese
desbalance es lo que vuelve tan efectivo al \emph{quishing}: el riesgo está ahí, pero
es invisible justo en el momento en que el usuario tiene que decidir. Por eso tiene
sentido atacar el problema donde ocurre, es decir, en la pantalla de escaneo.

El tamaño del problema y su evolución ya se documentaron en la definición del
problema, y el análisis de qué tan confiable es esa evidencia se desarrolla en los
antecedentes. Lo que toca argumentar aquí es otra cosa: por qué vale la pena
atenderlo ahora, por qué esta forma de resolverlo es la adecuada y a quién le sirve.
Sobre lo primero bastan tres datos ya establecidos. El ataque se industrializó al
integrarse a plataformas de \emph{phishing-as-a-service} \cite{insikt2024qrcode}. Los
filtros de correo corporativo no leen los códigos QR incrustados como imagen, así que
este vector simplemente los esquiva \cite{unit42_2023qr}. Y las apps de escaneo
seguro que existen funcionan casi todas con listas negras, un enfoque que deja fuera
a los dominios recién creados
\cite{sheng2009empirical,ma2009beyond,apwg2024q1,apwg2025q1}. Juntas, estas tres
condiciones dejan abierta una ventana de exposición que hoy nadie cubre.

\subsection*{Por qué intervenir en el escaneo y no en otro punto}

Hay una objeción obvia a esta propuesta: los navegadores móviles ya traen servicios
de reputación de URLs, entonces el usuario estaría protegido sin importar con qué app
escanee. Es una objeción válida y conviene responderla de frente, porque de ahí
depende que este trabajo tenga sentido.

El punto es que el navegador avisa tarde. Para cuando muestra la advertencia, la
petición al servidor del atacante ya salió: su dirección IP, el modelo de navegador y
todo lo que venía codificado en la URL ya se enviaron, y lo que el servidor haya
respondido ya llegó al teléfono. En una campaña de \emph{quishing}, buena parte de lo
que el atacante quiere lo consigue en ese primer contacto, sin que el usuario tenga
que hacer nada más. Si el análisis se mete entre la lectura del código y la apertura
del navegador, el control se adelanta un paso: ahí la URL ya se puede leer, pero
todavía no hubo contacto con el atacante. Esa diferencia de un solo paso es la que
separa una advertencia que llega tarde de una que llega a tiempo, y es la razón de
que la solución viva en la app de escaneo y no más adelante en la cadena.

De esa decisión se desprenden las demás. La lectura del código se resuelve en el
propio teléfono con las APIs nativas de visión
\cite{apple_vndetectbarcodes,google_mlkit_barcode,rousavy_visioncamera}, así no hay
que mandar la imagen a ningún lado y el tiempo total de respuesta se mantiene dentro
de lo aceptable. El análisis arranca consultando fuentes de reputación porque son las
que dan mayor certeza sobre amenazas ya conocidas gastando muy poco tiempo y cómputo
\cite{google_safebrowsing,virustotal_apiv3,phishtank_devinfo}; son el primer filtro
lógico. Que no alcancen para dominios nuevos es, justamente, la razón de que existan
filtros después. Y se consultan varias fuentes en lugar de una sola porque sus
coberturas son parciales y desparejas, sobre todo con campañas en español dirigidas a
la región, donde se reporta bastante menos que en inglés.

\subsection*{A quién le sirve este trabajo}

En lo teórico, el trabajo aporta al estudio de ataques que aprovechan interfaces poco
convencionales, donde la señal que el usuario normalmente usaría para desconfiar (ver
la URL) desaparece por diseño del medio. Ya está documentado que la gente cae más
fácil en phishing por QR que en phishing por enlace directo, precisamente por eso
\cite{vidas2013qrishing,krombholz2014qr}, lo que coloca el problema entre la
seguridad informática y el factor humano. Además es un terreno poco explorado en
América Latina, donde la tecnología se adoptó más rápido que las prácticas de
seguridad que deberían venir con ella.

Para los usuarios, una interfaz que revise el enlace entre la lectura del código y la
redirección atiende un riesgo que aparece en situaciones de todos los días:
comercios, transporte público, oficinas de gobierno y pagos. Como el quishing le pega
a cualquiera con un teléfono, sepa o no de tecnología \cite{vidas2013qrishing}, el
valor de una herramienta así está en que no le exige al usuario saber nada especial
para decidir bien.

Vale la pena aclarar hasta dónde llega esa afirmación. Lo que este trabajo aporta en
diseño es convertir un veredicto técnico en una decisión que el usuario pueda tomar,
en el momento y en el aparato donde le toca tomarla. Eso se mide con sus propios
criterios: si el veredicto se entiende, si el flujo funciona bien y si cumple el
tiempo de respuesta. Son criterios distintos a los que miden qué tan bien detecta el
motor de análisis. Las dos cosas hacen falta y ninguna reemplaza a la otra: una
pantalla muy clara que muestre un veredicto poco confiable no protege a nadie, y un
veredicto acertado pero mal explicado tampoco. Lo que se sostiene aquí es que a la
parte de comunicar se le ha puesto mucha menos atención que a la de detectar, y que
documentarla bien tiene valor propio. El desempeño real de las fuentes de reputación
se mide como parte de la evaluación y se reporta con sus limitaciones.

Para los negocios e instituciones que usan códigos QR para atender público, el
trabajo ofrece una forma documentada de entender un riesgo que golpea directo la
confianza en ese canal. Que alguien pegue una calcomanía encima de un QR legítimo
\cite{krombholz2014qr} es un problema que el establecimiento no puede resolver por su
cuenta; tener evidencia de cómo se puede verificar un enlace antes de abrirlo les da
una base para pensar qué medidas tomar.

Para la Universidad del Valle de Guatemala, este trabajo suma en un área,
la ciberseguridad aplicada, donde todavía hay poca producción académica local, y deja
un antecedente para futuros proyectos de seguridad móvil, análisis de amenazas y
educación digital. La falta de formación especializada en ciberseguridad es uno de
los factores que más limitan la capacidad de respuesta de los países en desarrollo
frente a amenazas nuevas; proyectos como este ayudan a cerrar esa brecha desde la
universidad, tanto por lo que documentan como por lo que enseñan a quienes los hacen.

Para quienes integramos el equipo, el proyecto significa afianzar conocimientos en
seguridad informática, desarrollo móvil multiplataforma e integración de servicios de
\emph{threat intelligence}, además de trabajar sobre un problema con implicaciones
éticas reales: qué se le advierte al usuario, con cuánto detalle y con cuánta certeza
son decisiones que afectan directamente su seguridad.

Para el público en general, el trabajo ayuda a dar visibilidad a una amenaza que ha
crecido bastante pero que todavía suena poco. Al describir el problema y construir
una solución que se puede evaluar con criterios claros, el proyecto deja material
útil para que instituciones públicas, empresas y centros educativos decidan con mejor
información en temas de seguridad digital.

En resumen, este trabajo se justifica porque coinciden tres cosas: una amenaza que
crece y que ya se volvió industria, un conjunto de soluciones existentes con
limitaciones técnicas bien documentadas, y una población expuesta que no tiene ninguna
protección en el momento que más importa. Atacar el problema desde la app móvil, el
único punto donde la URL se puede leer antes de que haya contacto con el atacante,
tiene sentido tanto por dónde está ubicado en la cadena del ataque como porque
permite evaluar la propuesta con criterios concretos de desempeño y usabilidad.
